\documentclass[10pt,a4paper]{amsart}
\usepackage[margin=19mm]{geometry}
\usepackage{amsmath,amssymb,mathtools}
\usepackage[scr=rsfs,cal=euler]{mathalfa}
\usepackage{xfrac}
\usepackage[unicode,hidelinks,bookmarks=false]{hyperref}
\newcommand{\ie}{i.e.\ }
\newcommand{\cf}{cf.\ }
\newcommand{\resp}{respectively\ }
\newcommand{\Subsection}{\subsection}
\providecommand{\nobreakdash}{\nobreak\hbox{-}}
\setlength{\emergencystretch}{2em}
\multlinegap0pt
\usepackage{amssymb}
\usepackage{xcolor}
\newtheorem{proposition}{Proposition}
\title{Uniformly affine actions on Banach spaces: growth of cocycles}
\author{Kevin Boucher, Georg Grützner}
\date{Source: arXiv:2601.06322v1 (CC BY 4.0)}
\begin{document}
\maketitle

\noindent\emph{Source and licence.} Uniformly affine actions on Banach spaces: growth of cocycles. Version arXiv:2601.06322v1 (CC BY 4.0), distributed by arXiv under the Creative Commons Attribution 4.0 International licence, http://creativecommons.org/licenses/by/4.0/, as recorded in the arXiv licence metadata for this exact version and shown on the abstract page https://arxiv.org/abs/2601.06322v1. Full licence notice: \texttt{licenses/arxiv-2601.06322-CC-BY-4.0.txt}.\par\smallskip

\noindent\textit{Editorial scope.} Counted unit: the article's proposition 2a47dc (source0.tex lines 941--948). The article's complete original proof of it is delivered with it (source0.tex lines 979--1030, one \texttt{proof} environment of 52 lines, which the article places away from the statement). No mathematics is written, repaired or re-derived: the statement and the proof are the article's own lines, in the article's own notation, and no retained line is substituted or re-flowed.  No further source-local context is needed: every cross-reference of every retained line resolves inside the two delivered spans.  Nothing was omitted from any retained span, so the package carries no omission record.  No retained line contains a citation, so the wrapper carries no bibliography and no bibliography entry.  No retained line contains a figure environment, an image-inclusion command or drawing code, so no image is delivered and the package declares no asset dependency.  Editorial formatting only, in addition: an AMS \texttt{amsart} preamble that restates the article's own declarations and macros used by the retained lines, namely the environments \texttt{proposition} on separate counters, together with the standard packages those lines need.  The retained environments print in this document as Proposition 1 in order of appearance, whereas the article prints the same items with the numbering of its own full document.  The complete native source of the article as distributed in the arXiv e-print for this exact version is bundled unchanged as \texttt{sources/192401/source0.tex}.
\begin{proposition}
\label{2a47dc}
Let \(F: G/K \to H\) be a nonconstant \(G\)-equivariant harmonic map.
Then the Hilbert-Schmidt norm of the operator
\(d_{x}F: T_{x}(G/K) \to H\) is uniformly bounded, i.e.~there exists
\(c > 0\) s.t. for all \(x \in G/K\),
\[ \frac{\sqrt{ \operatorname{dim}(G/K) }}{c} \leq \|d_{x}F\|_{HS} \leq c \, \sqrt{ \operatorname{dim}(G/K) }.\]
\end{proposition}

\begin{proof}[Proof of Proposition \ref{2a47dc}]
The tangent space \(T_{x_0}(G / K)\) is identified with \(\mathfrak{p}\)
as previously described, and the tangent space at any vector \(v \in H\)
is similarly identified with \(H\). Using the fact that \(F\) is
\(G\)-equivariant, we deduce the following for all \(k \in K\) and
\(X \in \mathfrak{p}\): \[
\begin{aligned}
d_{x_{0}} F(\operatorname{Ad}(k) X) & =\left.\frac{d}{d t} F\left(\exp(t \operatorname{Ad}(k) X) x_0\right)\right|_{t=0} \\
& =\left.\frac{d}{d t} F\left(k \exp(t X) x_0\right)\right|_{t=0} \\
& =\left.\frac{d}{d t}\left(\pi(k) F\left(\exp(t X) x_0\right)\right)\right|_{t=0} \\
& =\pi(k) \left.\frac{d}{d t} F\left(\exp(t X) x_0\right)\right|_{t=0} \\
& =\pi(k) \, d_{x_{0}} F(X).
\end{aligned}
\] Since \(K\) acts unitarily on \(H\),
\(\|\pi(k)d_{x_{0}} F(X)\| = \|d_{x_{0}} F(X)\|\) for all \(k \in K\).
Thus, the symmetric bilinear form \(B_{x_{0}}\) on
\(T_{x_0}(G / K) \cong \mathfrak{p}\), defined by \[
B_{x_{0}}(X, Y) = \left\langle d_{x_{0}} F(X), d_{x_{0}} F(Y) \right\rangle,
\] is invariant under the action of \(\operatorname{Ad}(K)\).

We claim that \(B_{x_{0}}\) is positive definite. Assume that there
exists a non-zero \(Z \in T_{x_{0}}(G/K)\) s.t.
\(\|d_{x_{0}}F(Z)\| = 0\). Then by the \(\operatorname{Ad}(K)\)
invariance of \(B_{x_{0}}\), \(d_{x_{0}}F\) vanishes on the whole
tangent space \(T_{x_{0}}(G/K)\). By \(G\)-equivariance of \(F\), for
any \(g \in G\) and \(x = g x_0\), it holds that \begin{equation}
\label{9e84cd}
d_{x} F = \pi(g) \,d_{x_{0}} F\,d_{x} \lambda\left(g^{-1}\right).\end{equation}

Since \((\pi,H)\) is uniformly bounded,
\[\left\|d_{x} F\left(d_{x_{0}} \lambda(g) (X)\right)\right\| \lesssim \left\|d_{x_{0}} F(X)\right\| = 0\]
for all \(X \in \mathfrak{p}\). Thus \(F\) would be constant, a
contradiction.

By norm equivalence, there exists \(c > 0\), s.t.
\[\frac{\|X\|_{2}^2}{c^{2}}\leq B_{x_{0}}(X,X) \leq c^{2} \,\|X\|_{2}^2 \quad \text{for all } X \in T_{x_0}(G / K).\]

We claim that \(c\) can be chosen independently of the base-point
\(x_{0}\). Indeed, by \(G\)-equivariance of \(F\), for any \(g \in G\)
and \(x = g x_0\), we have \[
\begin{aligned}
B_{x}(d_{x_{0}} \lambda(g) (X),d_{x_{0}} \lambda(g) (X) )
&=\left\|d_{x} F\left(d_{x_{0}} \lambda(g) (X)\right)\right\|^{2} \\
& \leq C \, \left\|d_{x_{0}} F(X)\right\|^{2} \\
& = C \, B_{x_{0}}(X,X) \\
& \leq C\,c^{2} \, \left\|X\right\|_2^{2}
\end{aligned}
\] for all \(X \in \mathfrak{p}\). A lower bound follows analogously.
The uniform bounds on the Hilbert-Schmidt norm follows by summing over
an orthonormal basis.

\end{proof}

\end{document}
